\documentclass[a4paper]{article}
\usepackage[pdftex]{graphicx}
\usepackage[utf8]{inputenc}
\usepackage{enumerate}
\usepackage{siunitx}
\sisetup{locale=DE} 
\usepackage{href-ul}
\hypersetup{
	colorlinks=true,
	linkcolor=blue,
	urlcolor=blue}
\usepackage{icomma}
\usepackage{amssymb}
\usepackage{geometry}
\geometry{a4paper, top=15mm, left=15mm, right=15mm, bottom=15mm,
	headsep=10mm, footskip=12mm}
%Source: img_0240

% Start the document
\begin{document}
%Titel
{\bf Elektrische Feldstärke und das Coulomb'sche Gesetz }\\
\begin{enumerate}[1.]

%Aufgabe 1
\item {\bf Wasserstoffatom}
\begin{enumerate}[a)]
\item Berechnen Sie den Betrag der elektrischen Kraft $F_C$ zwischen Atomkern und Elektron eines Wasserstoffatoms. Für den Abstand $r$ des Elektrons vom Proton gilt: $r=\SI{5,3e-11}{\meter}$ (=Radius der Bohr-Grundbahn).
\item Vergleichen Sie diesen Wert $F_C$ mit der Gravitationskraft $F_G$ zwischen Atomkern und Elektron.
\end{enumerate}

%Aufgabe 2
\item {\bf Aufhebung von Gravitationskraft und Coulombkraft}\\
Zwei als Massepunkte zu betrachtende Körper der Masse $m_1$ und $m_2$ tragen die positiven Ladungen $Q_1$ und $Q_2$ mit $m_1=m_2=\SI{1,0e-3}{\kilogram} $ und \\
$Q_1= \SI{1,0e-13}{\coulomb}$. Wie groß muss die Ladung $Q_2$ sein, damit sich Gravitationskraft und Coulombkraft aufheben?

%Aufgabe 3
\item {\bf Pendel mit zwei gleichnamig geladenen Kugeln}\\
Zwei gleiche Kugeln mit der Gewichtskraft von je $F_G=\SI{0,5e-2}{\newton}$ sind an zwei je 
$l=\SI{1,00}{\meter}$ langen, oben an demselben Punkt befestigten Fäden aufgehängt und tragen gleiche Ladungen $Q_1=Q_2=Q$. Die Kugeln haben wegen der Abstoßung den Abstand $d=\SI{0,40}{\meter}$. Berechnen Sie den Betrag der Ladung auf den Kugeln.

%Aufgabe 4
\item {\bf Positiv geladene Metallkugel}\\
Eine Metallkugel vom Radius $R=\SI{4,5}{\centi\meter}$ trägt eine positive Ladung \\ 
$Q=\SI{1,0e-9}{\coulomb}$.
\begin{enumerate}[a)]
\item Geben Sie den Betrag der elektrischen Feldstärke mit eingesetzten Zahlenwerten in Abhängigkeit vom Abstand $r$ an.
\item Berechnen Sie den Betrag der elektrischen Feldstärke im Bereich \\ $R \le r \le 3 R$.
\item Stellen Sie den Verlauf der elektrischen Feldstärke für $R \le r \le 3 R$ graphisch dar.
\item Geben Sie den Betrag der elektrischen Feldstärke im Inneren der Kugel an und ergänzen Sie entsprechend das Diagramm von Teilaufgabe c).
\item Berechnen Sie die Kraft auf einen Heliumkern ($\alpha$-Teilchen), der sich in der Entfernung $r_{\alpha}=\SI{13,5}{\centi\meter}$ vom Kugelmittelpunkt befindet.
\end{enumerate}

%Aufgabe 5
\item {\bf Luftballon}\\
ein kugelförmiger Luftballon ist außen mit Graphit leitend bestrichen und isoliert aufgehängt. Sein Radius $r_1$ beträgt $\SI{2,5}{\centi\meter}$. Auf dem Ballon ist die Ladung $Q$ gleichmäßig verteilt. In der Entfernung $d=\SI{52,5}{\centi\meter}$ vom Kugelmittelpunkt ruft diese Ladung die elektrische Feldstärke $E_1=\SI{100}{\newton\over \coulomb}$ hervor.
 \begin{enumerate}[a)]
\item Berechnen Sie den Betrag der Ladung $Q$.
\item Der Ballon wird nun aufgeblasen, der neue Radius $r_2$ beträgt $\SI{5,0}{\centi\meter}$. Ermitteln Sie die Feldstärke $E_2$, die sich in der Entfernung $d$ vom Kugelmittelpunkt einstellt.
\end{enumerate} 

%Aufgabe 6
\item {\bf Feldstärkefreier Punkt}\\
Gegeben sind zwei positive Ladungen mit $Q_1=\SI{8,5}{\nano\coulomb}$ und $Q_2=\SI{5,5}{\nano\coulomb}$. Auf der Verbindungsgeraden der beiden Ladungen befindet sich im Abstand $s=\SI{10,0}{\centi\meter}$ von der Ladung $Q_1$ ein Punkt $P$, in dem die elektrische Feldstärke Null ist. Berechnen Sie den Abstand $d$ der beiden Punktladungen.
\end{enumerate} 

\fbox{
	\begin{minipage}{0.5\textwidth}
		Zur Lösung bitte \href{https://www.okuyakl.de/phys/p11coulL240/ll240.pdf}{hier klicken} oder den QR-Code scannen.\\
	Weitere Arbeitsblätter gibt es unter 
	
	\href{https://www.okuyakl.de}{www.okuyakl.de}
	\end{minipage}
	\hfill
	\begin{minipage}{0.4\textwidth}
		\includegraphics[width=1.5 cm]{../../viecher/zwe03}
		\includegraphics[width=3 cm]{qr240}
		\includegraphics[width=2 cm]{../../viecher/afanticon1}
		
	\end{minipage}}

\end{document}%Lösung--------------------------------------------
